% Options for packages loaded elsewhere
\PassOptionsToPackage{unicode}{hyperref}
\PassOptionsToPackage{hyphens}{url}
\PassOptionsToPackage{dvipsnames,svgnames,x11names}{xcolor}
%
\documentclass[
  11pt,
]{article}
\usepackage{amsmath,amssymb}
\usepackage{iftex}
\ifPDFTeX
  \usepackage[T1]{fontenc}
  \usepackage[utf8]{inputenc}
  \usepackage{textcomp} % provide euro and other symbols
\else % if luatex or xetex
  \usepackage{unicode-math} % this also loads fontspec
  \defaultfontfeatures{Scale=MatchLowercase}
  \defaultfontfeatures[\rmfamily]{Ligatures=TeX,Scale=1}
\fi
\usepackage{lmodern}
\ifPDFTeX\else
  % xetex/luatex font selection
\fi
% Use upquote if available, for straight quotes in verbatim environments
\IfFileExists{upquote.sty}{\usepackage{upquote}}{}
\IfFileExists{microtype.sty}{% use microtype if available
  \usepackage[]{microtype}
  \UseMicrotypeSet[protrusion]{basicmath} % disable protrusion for tt fonts
}{}
\makeatletter
\@ifundefined{KOMAClassName}{% if non-KOMA class
  \IfFileExists{parskip.sty}{%
    \usepackage{parskip}
  }{% else
    \setlength{\parindent}{0pt}
    \setlength{\parskip}{6pt plus 2pt minus 1pt}}
}{% if KOMA class
  \KOMAoptions{parskip=half}}
\makeatother
\usepackage{xcolor}
\usepackage[margin=0.9in]{geometry}
\usepackage{longtable,booktabs,array}
\usepackage{calc} % for calculating minipage widths
% Correct order of tables after \paragraph or \subparagraph
\usepackage{etoolbox}
\makeatletter
\patchcmd\longtable{\par}{\if@noskipsec\mbox{}\fi\par}{}{}
\makeatother
\setlength{\emergencystretch}{3em} % prevent overfull lines
\providecommand{\tightlist}{%
  \setlength{\itemsep}{0pt}\setlength{\parskip}{0pt}}
\setcounter{secnumdepth}{-\maxdimen} % remove section numbering
\ifLuaTeX
  \usepackage{selnolig}  % disable illegal ligatures
\fi
\usepackage{hyperref}
\IfFileExists{xurl.sty}{\usepackage{xurl}}{} % add URL line breaks if available
\urlstyle{same}
\hypersetup{
  colorlinks=true,
  linkcolor={Maroon},
  filecolor={Maroon},
  citecolor={Blue},
  urlcolor={Blue},
  pdftitle={Exact degrees for even Catalan-power Hankel determinants},
  pdfauthor={Aurantium},
  pdfcreator={LaTeX}}

\title{Exact degrees for even Catalan-power\\Hankel determinants}
\author{Aurantium}
\date{13 October 2026}

\begin{document}
\maketitle

\begin{abstract}
For integers \(j\ge1\) and \(k\ge2j-1\), the signed Hankel determinant
sequences formed from the coefficients of the even Catalan power
\(C(z)^{2k}\), with zero moment shift and matrix sizes \(kn+j\) and
\(kn+k+1-j\), agree for every \(n\ge0\) with rational polynomials of
exact degree \((2j-1)(k-j)\). The sign is
\((-1)^{n\binom{k}{2}}\). A quotient-minor formula gives the degree bound;
a strictly positive moment determinant proves that its leading term does
not cancel. We give the full proof, including the boundary cases, and
exact finite verification code. The result concerns these two offsets
in the stated parameter range, not arbitrary moment shifts, odd Catalan
powers, or Narayana-parameter deformations.
\end{abstract}

\section{Statement and notation}\label{problem}

\subsection{The residue-three case}\label{exact-reconstruction}

Let \[
 C(z)=\frac{1-\sqrt{1-4z}}{2z},\qquad
 c_m^{(2k)}=[z^m]C(z)^{2k}
 =\frac{2k}{2m+2k}\binom{2m+2k}{m},
\] and \[
 D_{2k,0}(N)=\det(c_{u+v}^{(2k)})_{0\le u,v<N},\qquad
 D_{2k,0}(0)=1.
\] For integers \(k\ge5\), \(n\ge0\), and \(r\in\{3,k-2\}\), set \[
 f_{k,r}(n)=(-1)^{nK}D_{2k,0}(kn+r),\qquad K=\binom{k}{2}.
\] The target is the existence of polynomials \(F_{k,r}\in\mathbb Q[n]\)
such that \[
 F_{k,r}(n)=f_{k,r}(n)\quad(n\ge0),\qquad
 \deg F_{k,r}=5(k-3).
\] At \(k=5\), the two residues are identical. The asserted degree is
exact, not merely an upper bound. The polynomials are integer-valued on
the indicated integer arguments; integer coefficients are not asserted.

This is precisely the \(j=3\) specialization of {[}C23, Conjecture 6,
equation (19){]}, with the sign convention in equation (16). The
specializations with \(5\le k\le15\) are
already included in {[}WX19, Theorem 11{]}.

\section{Theorems and proofs}\label{result}

\subsection{1. Resolution and leading
coefficient}\label{resolution-and-leading-coefficient}

\textbf{Theorem 1.} The target holds for every \(k\ge5\) and every
\(n\ge0\). Define \[
 b_m=\frac{2\zeta(2m)}{\pi^{2m}}
     =\frac{(-1)^{m+1}2^{2m}B_{2m}}{(2m)!}>0\qquad(m\ge1),
 \tag{1}
\] where \(u/(e^u-1)=\sum_{m\ge0}B_mu^m/m!\). The leading coefficients
are \[
 \boxed{\ [n^{5(k-3)}]F_{k,3}(n)
 =L_k:=k^{4(k-3)}\bigl(b_{k-4}b_{k-2}-b_{k-3}^{\,2}\bigr)>0,\ }
 \tag{2}
\] \[
 \boxed{\ [n^{5(k-3)}]F_{k,k-2}(n)=(-1)^K L_k.\ }
 \tag{3}
\] In particular, there is no unproved noncancellation assumption in the
degree claim. Strict positivity in (2) is established below.

The proof also extends to every residue pair in the stated range of
{[}C23, Conjecture 6{]}. Section 7 gives a complete proof of that
stronger conclusion, not just a proposed extension. No claim of
publication priority is made.

\subsection{2. Moment functional and a finite determinant
reduction}\label{moment-functional-and-a-finite-determinant-reduction}

Put \(s=k-1\). Define monic polynomials by \[
 P_0(x)=1,\quad P_1(x)=x-2,\quad
 P_{m+1}(x)=(x-2)P_m(x)-P_{m-1}(x).
 \tag{4}
\]

\textbf{Lemma 2.1 (moment bridge).} There is a linear functional
\(\mathcal L\) on \(\mathbb Q[x]\) such that \[
 \mathcal L(P_aP_b)=\delta_{ab},\qquad
 c_m^{(2k)}=\mathcal L\bigl(x^m Q(x)\bigr),\qquad
 Q(x)=x^sP_s(x).
 \tag{5}
\] Moreover, \(H_N(\mathcal L):=\det(\mathcal L(x^{u+v}))_{u,v<N}=1\).

\emph{Proof.} Make the \(P_m\) an orthonormal basis of the polynomial
space. The recurrence (4) says that multiplication by \(x\) is
symmetric: its matrix has diagonal entries 2 and adjacent off-diagonal
entries 1. At the boundary, \(xP_0=P_1+2P_0\). Define
\(\mathcal L(h)=\langle1,h\rangle\). Symmetry of multiplication gives
\(\langle f,g\rangle=\mathcal L(fg)\), proving orthogonality.

The coefficient of \(P_a\) in \(x^m\) counts nonnegative Motzkin walks
of length \(m\) from height 0 to height \(a\), with horizontal-step
weight 2 and both sloping-step weights 1. Let \(G(z)\) be the excursion
generating function. Decomposition at the first step and first return
gives \[
 G=1+2zG+z^2G^2.
\] This has a unique formal solution with constant coefficient 1. Since
\(C=1+zC^2\), its solution is \(G=C^2\).

Decomposing a walk ending at height \(s\) at its last upward departures
from heights \(0,\ldots,s-1\) gives \(s\) up-steps and \(s+1\)
excursions. Thus \[
 \sum_{m\ge0}\mathcal L(x^mP_s)z^m=z^sG(z)^{s+1}.
\] Taking the coefficient of \(z^{m+s}\) proves (5), because \(s+1=k\).
Finally, the change of basis from monomials to the monic \(P_m\) is unit
triangular. Their Gram matrix is the identity, so \(H_N(\mathcal L)=1\),
including the empty determinant. \(\square\)

For any integer \(N\ge0\), let \(\mathcal R(N)\) be the determinant of
the coefficient vectors of \[
 P_N,P_{N+1},\ldots,P_{N+2s-1}\pmod Q
\] in the ordered monomial basis \(1,x,\ldots,x^{2s-1}\).

\textbf{Lemma 2.2 (remainder determinant).} For every \(N\ge0\), \[
 D_{2k,0}(N)=\mathcal R(N).
 \tag{6}
\]

\emph{Proof.} Consider the \((N+2s)\)-column coefficient matrix of \[
 QP_0,\ldots,QP_{N-1},P_N,\ldots,P_{N+2s-1}.
\] In the orthonormal \(P\)-basis, its last \(2s\) columns are the
corresponding unit vectors. Its determinant is therefore
\(\det(\mathcal L(P_uQP_v))_{u,v<N}=H_N(Q\mathcal L)\). Changing the
first \(N\) monic basis polynomials to monomials does not change this
determinant. By Lemma 2.1 it is \(D_{2k,0}(N)\).

For each last column, division by \(Q\) has quotient degree at most
\(N-1\). Subtract a linear combination of the first \(N\) columns to
leave its remainder. Move these \(2s\) remainder columns to the front.
The sign is \((-1)^{2sN}=1\). The resulting matrix has remainder block
in its upper left, zero lower-left block, and a unit-triangular
lower-right block. Its determinant is \(\mathcal R(N)\). No nonvanishing
of \(H_N(Q\mathcal L)\) was assumed. \(\square\)

This reduction is a specialization of the classical
orthogonal-polynomial moment-modification identity in {[}K21, Theorem 1
and Proposition 5{]}. The proof above is included to check its
hypotheses and conventions directly; the reduction itself is not
asserted to be new.

\subsection{3. Evaluation matrix, polynomiality, and
reflection}\label{evaluation-matrix-polynomiality-and-reflection}

The roots of \(P_s\) are \[
 \alpha_i=2+2\cos(i\pi/k),\qquad 1\le i\le s.
\] They are distinct and strictly positive. In particular they are
different from zero. Let \[
 V(\alpha)=\prod_{1\le i<j\le s}(\alpha_j-\alpha_i).
\] For a polynomial of degree below \(2s\), take its values at the \(s\)
roots, followed by its coefficients of degrees \(0,\ldots,s-1\). The
determinant of this linear transformation on the monomial basis is \[
 V(\alpha)P_s(0)^s\ne0.
 \tag{7}
\] Indeed, move the \(s\) coefficient rows to the front. Their block is
\((I_s\ 0)\); the remaining block in degrees \(s,\ldots,2s-1\) has
determinant \((\prod_i\alpha_i^s)V(\alpha)\). The row-move sign
\((-1)^{s^2}\) gives exactly (7), since
\(P_s(0)=(-1)^s\prod_i\alpha_i\).

A multiple of \(Q=x^sP_s\) is annihilated by all these functionals.
Consequently, if \(E_N\) is the matrix with columns
\(P_N,\ldots,P_{N+2s-1}\), root-value rows first and coefficient rows
second, then \[
 \mathcal R(N)=\frac{\det E_N}{V(\alpha)P_s(0)^s}.
 \tag{8}
\]

Extend (4) to all integer indices. It gives \[
 P_{-1}=0,\qquad P_{-m-2}=-P_m\quad(m\ge0).
 \tag{9}
\] Use the remainder determinant to define \(\mathcal R(N)\) also for
negative \(N\). This is an algebraic extension, not a negative-size
Hankel determinant. Two useful identities, valid for integer \(m\), are
\[
 P_m(\alpha_i)=\frac{\sin((m+1)i\pi/k)}{\sin(i\pi/k)},\qquad
 [x^a]P_m(x)=(-1)^{m-a}\binom{m+a+1}{2a+1}.
 \tag{10}
\] The binomial in the second identity is the falling-factorial
polynomial. For nonnegative \(m\), the identities follow from (4) and
its initial conditions. For negative \(m\), they follow from (9) and
\(\binom{-u}{v}=(-1)^v\binom{u+v-1}{v}\); the index \(-1\) gives zero.

\textbf{Lemma 3.1.} For any fixed integer residue \(r\), \[
 (-1)^{nK}\mathcal R(kn+r)\quad(n\in\mathbb Z)
\] is the restriction of a polynomial \(F_{k,r}\in\mathbb Q[n]\) of
degree at most \(s^2\).

\emph{Proof.} By (10), root row \(i\) of \(E_{kn+r}\) has a common
factor \((-1)^{ni}\). Their product is \((-1)^{nK}\). Removing it makes
the root rows constant in \(n\).

Coefficient row \(a\), \(0\le a<s\), has a common sign \((-1)^{kn}\);
after removing it, each entry is a polynomial of degree at most
\(2a+1\). The product of these signs is 1, because \(kns=k(k-1)n\) is
even. The determinant in (8) thus has degree at most
\(\sum_{a=0}^{s-1}(2a+1)=s^2\). It agrees with the signed remainder
determinant at every integer \(n\), including negative integers. Its
coefficients are rational: its values at any \(s^2+1\) distinct integers
are integers, and the corresponding Vandermonde interpolation system is
over \(\mathbb Q\). \(\square\)

\textbf{Lemma 3.2 (reflection).} As polynomial identities, \[
 F_{k,k+1-r}(n)=(-1)^{s+K}F_{k,r}(-n-3).
 \tag{11}
\]

\emph{Proof.} Reverse the \(2s\) columns and apply (9). The \(2s\)
individual minus signs cancel, and reversal has sign \((-1)^s\).
Therefore \[
 \mathcal R(-N-2s-1)=(-1)^s\mathcal R(N).
\] If \(N=kn+r\), its reflected index is \[
 -N-2s-1=k(-n-3)+(k+1-r).
\] Combining the two signing factors gives \((-1)^{s+K}\). Lemma 3.1
promotes the identity on integer arguments to the polynomial identity
(11). \(\square\)

\subsection{4. The residue-3 quotient
identity}\label{the-residue-3-quotient-identity}

For \(n\ge0\), write \[
 N=kn+3,\qquad M=k(n+1),\qquad
 p(x)=P_{M-1}(x),\quad q(x)=P_{M+k-1}(x),
\] \[
 \frac{q(x)}{p(x)}=\sum_{a\ge0}h_a(k,M)x^a.
 \tag{12}
\] The quotient is a well-defined formal series because
\(p(0)=(-1)^{M-1}M\ne0\).

\textbf{Lemma 4.1.} For every \(k\ge5\) and \(n\ge0\), \[
 \boxed{\ F_{k,3}(n)=-(n+1)^s
 \left(h_{s-2}(k,M)^2-h_{s-3}(k,M)h_{s-1}(k,M)\right).\ }
 \tag{13}
\]

\emph{Proof.} Let \(\ell=M-1=kn+s\). In the column window
\(P_N,\ldots,P_{N+2s-1}\), replace the columns with leading indices
\(\ell,\ldots,\ell+s-3\) by \[
 p,xp,\ldots,x^{s-3}p.
\] These are determinant-preserving triangular column operations.
Specifically,
\(x^aP_\ell=P_{\ell+a}+\text{a combination of }P_{\ell-a},\ldots,P_{\ell+a-1}\),
and \(\ell-a\ge N\) for \(a\le s-3\).

The final two original columns are \(q=P_{kn+2s+1}\) and
\(P_{kn+2s+2}\). Replace the latter by \[
 xq=P_{kn+2s+2}+2P_{kn+2s+1}+P_{kn+2s}.
\] This is again determinant-preserving. Both \(p\) and \(q\) vanish at
every \(\alpha_i\), because their indices plus one are multiples of
\(k\).

The \(s\) columns which remain nonzero in root rows are \[
 P_{kn+3},\ldots,P_{kn+s-1},\quad
 P_{kn+2s-2},P_{kn+2s-1},P_{kn+2s}.
\] After removing row factors \((-1)^{ni}\), these become \[
 P_3,\ldots,P_{s-1},-P_2,-P_1,-P_0.
\] Their root determinant is \((-1)^{s-3}V(\alpha)\). To see the sign
explicitly: there are three minus signs, \(3(s-3)\) crossings to move
the last three polynomials to the front, and three crossings to reverse
their order.

Moving the three right-hand root columns across the \(s-2\) columns made
from \(p\) has sign \((-1)^{3(s-2)}\). Multiplying these two signs gives
\((-1)^{s-3+3(s-2)}=-1\). After all root columns have been put first,
the lower-right coefficient block has columns \[
 p,xp,\ldots,x^{s-3}p,\ q,xq.
\]

Multiply the truncated coefficient vectors by \(1/p\), modulo \(x^s\).
The determinant of this multiplication is \(p(0)^{-s}\), since its
matrix is triangular. The first \(s-2\) columns become the corresponding
coordinate vectors, leaving the minor \[
 \begin{vmatrix}
 h_{s-2}&h_{s-3}\\
 h_{s-1}&h_{s-2}
 \end{vmatrix}.
\] The denominator in (8) supplies \(P_s(0)^s\), and \[
 \left(\frac{p(0)}{P_s(0)}\right)^s
 =\left((-1)^{kn}(n+1)\right)^s=(n+1)^s.
\] The root-row sign is exactly the one already removed in the
definition of \(F_{k,3}\). This proves (13), including \(n=0\).
\(\square\)

\subsection{5. Exact degree and strict
noncancellation}\label{exact-degree-and-strict-noncancellation}

The classical partial fraction expansion of cotangent {[}DLMF, 4.22.3{]}
is \[
 \cot u=\frac1u+2u\sum_{\nu\ge1}\frac1{u^2-\pi^2\nu^2}.
\] Expanding for \(0<|u|<\pi\) gives \[
 \cot u=\frac1u-\sum_{a\ge1}b_a u^{2a-1},\qquad
 b_a=2\sum_{\nu\ge1}(\pi\nu)^{-2a}.
 \tag{14}
\] The expansion and interchange of sums are justified by absolute
uniform convergence on compact subdisks of radius below \(\pi\), after
removal of the simple pole. The Bernoulli expression (1) is the
equivalent Laurent coefficient formula {[}DLMF, 4.19.6{]}.

\textbf{Lemma 5.1.} For fixed \(k\) and every \(a\ge1\), \(h_a(k,M)\) is
a Laurent polynomial in \(M\) with no exponent below \(-1\), with degree
\(2a-1\), and \[
 [M^{2a-1}]h_a(k,M)=(-1)^{k+1}k b_a.
 \tag{15}
\]

\emph{Proof.} Make the formal substitution \[
 x=2-2\cos\theta=\theta^2+O(\theta^4).
\] More precisely, \(\theta^2\) is the unique series in
\(x\mathbb Q[[x]]\) inverting this substitution, and has leading
coefficient 1. From (4), \[
 P_m(2-2\cos\theta)=(-1)^m
       \frac{\sin((m+1)\theta)}{\sin\theta}.
\] Consequently (12) gives \[
 \frac{q}{p}=(-1)^k\frac{\sin((M+k)\theta)}{\sin(M\theta)}
 =(-1)^k\left(\cos(k\theta)+\sin(k\theta)\cot(M\theta)\right).
 \tag{16}
\] This identity holds for the positive integer values of \(M\) used in
the Hankel determinant, and its right side defines the required
coefficient extension in the formal ring \(\mathbb Q[M,M^{-1}][[x]]\).

In (14), the pole contributes only a term proportional to \(M^{-1}\).
The cosine contributes a term independent of \(M\). For the coefficient
of \(x^a\), only cotangent terms indexed by \(1\le v\le a\) can
contribute, because \(\sin(k\theta)\theta^{2v-1}\) starts with
\(k\theta^{2v}\). Thus the possible exponents of \(M\) are
\(-1,0,1,3,\ldots,2a-1\). The highest one comes only from \(v=a\), and
\[
 [x^a]\bigl(\sin(k\theta)\theta^{2a-1}\bigr)=k.
\] This proves (15), with a nonzero coefficient. \(\square\)

Now put \(t=k-3=s-2\). The hypothesis \(k\ge5\) is exactly what ensures
\(t-1\ge1\), so Lemma 5.1 applies to all three coefficients in (13).
Their products have no pole worse than \(M^{-2}\), and have degree at
most \(4t-2\). Multiplication by \((n+1)^s=(M/k)^s\) clears all possible
poles. The degree is at most \[
 s+4t-2=5(k-3).
\] By (15), the coefficient at that degree is \[
 -k^{4t}\bigl(b_t^2-b_{t-1}b_{t+1}\bigr)=L_k.
\]

It remains to prove that this coefficient does not vanish. Put
\(u_\nu=(\pi\nu)^{-2}\). Because \(t\ge2\), all sums below converge
absolutely, and (14) gives the exact identity \[
 b_{t-1}b_{t+1}-b_t^2
 =4\sum_{1\le a<b}u_a^{t-1}u_b^{t-1}(u_a-u_b)^2>0.
 \tag{17}
\] Every term is nonnegative, and already the term \((a,b)=(1,2)\) is
positive. This proves (2) and proves that the degree is exact. No finite
computation is needed for this strict inequality.

Finally, apply (11) with \(r=3\). Reflection preserves degree, and the
leading coefficient changes by \[
 (-1)^{s+K+5(k-3)}=(-1)^K,
\] because \(s+5(k-3)=6k-16\) is even. This proves (3) and completes
Theorem 1. \(\square\)

\subsection{6. Boundary tests and exact
verification}\label{boundary-tests-and-exact-verification}

At the smallest permitted parameter \(k=5\), the two residues coincide.
The quotient formula gives the following exact polynomial: \[
 \begin{split}
 F_{5,3}(n)={}&\frac{125}{3024}(n+1)^2(n+2)^2(2n+3)^2\\
 &\quad\cdot\bigl(500n^4+3000n^3+5125n^2+1875n-1092\bigr).
 \end{split}
 \tag{18}
\] It has degree 10, leading coefficient \(15625/189\), and value
\(-1625\) at \(n=0\), equal to the original \(3\times3\) determinant.
Its reflection symmetry is \(F_{5,3}(n)=F_{5,3}(-n-3)\), as required.

The lower bound on \(k\) must not be dropped. The excluded pair
\(k=4,j=3\) has degree 6, not \(5(k-3)=5\); it is the reflected residue
of \(j=2\). This agrees with the printed small-parameter formulas in
{[}WX19, Theorem 6{]} and {[}C18, page 30{]}. It is not a counterexample
to Theorem 1, whose range is \(k\ge5\).

The accompanying standard-library program \texttt{verify\_cigler.py} performs
exact integer, rational, and Laurent-polynomial arithmetic. It does not
use floating point, a numerical zeta evaluator, or polynomial fitting.
Its deterministic data files contain the following recorded checks.
The complete suite was not rerun when preparing this research package;
the counts below describe the included records and reproducible code.

\begin{longtable}[]{@{}
  >{\raggedright\arraybackslash}p{(\columnwidth - 2\tabcolsep) * \real{0.5000}}
  >{\raggedright\arraybackslash}p{(\columnwidth - 2\tabcolsep) * \real{0.5000}}@{}}
\toprule\noalign{}
\begin{minipage}[b]{\linewidth}\raggedright
Check
\end{minipage} & \begin{minipage}[b]{\linewidth}\raggedright
Recorded range / count
\end{minipage} \\
\midrule\noalign{}
\endhead
\bottomrule\noalign{}
\endlastfoot
Symbolic residue-3 polynomials & 23 values of \(k\): \(5\) through
\(24\), and \(32,40,50\) \\
Original Hankel determinant comparisons & 60, for \(5\le k\le12\),
\(0\le n\le3\), both distinct target residues \\
Independent remainder-determinant comparisons & 429, for
\(5\le k\le24\), \(-4\le n\le6\), both distinct target residues \\
General-\(j\) quotient comparisons & 267, for \(2\le k\le18\), every
\(1\le j\le\lfloor(k+1)/2\rfloor\), \(0\le n\le2\) \\
Supplementary \(k=16\) certificate & All 452 values and both Newton
coefficient lists reproduced \\
Empty Hankel determinant & Checked to equal 1 \\
\end{longtable}

The negative-index checks use the algebraic remainder determinant, never
a negative-size Hankel determinant. They include \(n=-1\), where the raw
quotient has a zero denominator but its polynomial extension is
well-defined. The symbolic check constructs the coefficients of
\(\binom{M+a}{2a+1}\) as polynomial products and carries out formal
series division, independently of the trigonometric derivation of their
top degree.

\subsection{7. Stronger theorem: the full stated residue-degree
range}\label{stronger-theorem-the-full-stated-residue-degree-range}

The same construction proves the following generalization of the \(j=3\)
case.

\textbf{Theorem 7.1.} Let \(j\ge1\) and \(k\ge2j-1\) be integers. For
every \(n\ge0\), the two signed sequences in residues \(j\) and
\(k+1-j\) agree with polynomials of exact degree \[
 (2j-1)(k-j).
 \tag{19}
\] For \(k\ge2\), set \[
 d=j-1,\qquad t=k-j,\qquad m_0=k-2j+2.
\] If \(d\ge1\), the leading coefficient of \(F_{k,j}\) is \[
 \boxed{\ (-1)^{kd}k^{2dt}
 \det\bigl(b_{m_0+a+b}\bigr)_{0\le a,b<d}.\ }
 \tag{20}
\] The determinant in (20) is strictly positive. For \(d=0\), the
determinant is interpreted as 1, and the leading coefficient is 1.

\emph{Proof.} The case \(k=j=1\) is immediate from the excursion moments
in Lemma 2.1: \(D_{2,0}(N)=1\). Suppose henceforth that \(k\ge2\).

Take \(N=kn+j\), \(\ell=kn+s=M-1\), and again use \(p=P_\ell\) and
\(q=P_{kn+2s+1}\). Replace the columns starting at \(P_\ell\) by \[
 p,xp,\ldots,x^{s-j}p
\] and the columns starting at \(q\) by \[
 q,xq,\ldots,x^{j-2}q.
\] For \(j=1\), the second list is empty. Each replacement is a
unit-triangular column operation: the smallest required index in the
first list is \(\ell-(s-j)=N\); in the second list it is
\(kn+2s+1-(j-2)=kn+2s-j+3\), which is in the remaining right-hand
window. The upper bounds on required indices are exactly the columns
being replaced. Thus no polynomial outside the original window is
introduced.

After removal of root-row signs, the surviving \(s\) root columns are \[
 P_j,\ldots,P_{s-1},-P_{j-1},\ldots,-P_0.
\] Their determinant relative to \(V(\alpha)\) has sign \[
 (-1)^{j+j(s-j)+\binom j2}.
\] Moving the rightmost \(j\) root columns across the \(s-j+1=s-d\)
columns made from \(p\) has sign \((-1)^{j(s-j+1)}\). The combined sign
is therefore \((-1)^{\binom j2}\).

Divide the lower coefficient block by \(p\), modulo \(x^s\), just as in
Lemma 4.1. Its first \(s-d\) columns become coordinate vectors, while
its last \(d\) columns become \(h,xh,\ldots,x^{d-1}h\), where \(h=q/p\).
Thus the exact quotient identity is \[
 F_{k,j}(n)=(-1)^{\binom j2}(n+1)^s
 \det\bigl(h_{t+a-b}(k,M)\bigr)_{0\le a,b<d}.
 \tag{21}
\] For \(d=0\), it gives \(F_{k,1}(n)=(n+1)^s\), with the asserted exact
degree.

For \(d\ge1\), the least coefficient index in (21) is \[
 t-(d-1)=m_0=k-2j+2\ge1.
\] Therefore Lemma 5.1 applies to every entry. Each determinant term has
a possible pole of order at most \(d\), cleared by \(M^s\) because
\(s\ge d\). Its degree in \(M\) is at most \[
 \sum_{a=0}^{d-1}\bigl(2(t+a-\sigma(a))-1\bigr)=d(2t-1)
\] for every permutation \(\sigma\). Since \(s=t+d\), the total degree
is at most \[
 s+d(2t-1)=(2d+1)t=(2j-1)(k-j).
\]

The coefficient at this degree in the inner determinant is \[
 (-1)^{d(k+1)}k^d\det(b_{t+a-b})_{0\le a,b<d}.
\] Reversing its columns gives \[
 \det(b_{t+a-b})=(-1)^{\binom d2}
                  \det(b_{m_0+a+b}).
\] Combining this sign with \((-1)^{\binom{d+1}{2}}\) from (21), and
replacing \(M\) by \(k(n+1)\), gives exactly (20). The sign
simplification is \[
 \binom{d+1}{2}+d(k+1)+\binom d2\equiv kd\pmod2.
\]

To prove strict positivity, let \(v=(v_0,\ldots,v_{d-1})\in\mathbb R^d\)
be nonzero. Formula (14) yields \[
 \sum_{a,b=0}^{d-1}v_av_b b_{m_0+a+b}
 =2\sum_{\nu\ge1}u_\nu^{m_0}
      \left(\sum_{a=0}^{d-1}v_a u_\nu^a\right)^2>0.
 \tag{22}
\] The series converges since \(m_0\ge1\) and the polynomial in
parentheses is bounded on the bounded set of points \(u_\nu\). A nonzero
polynomial has only finitely many real zeros, whereas the \(u_\nu\) are
infinitely many distinct positive numbers. Hence at least one summand is
positive. This proves that the matrix in (20) is positive definite, so
its determinant is strictly positive. The degree bound is consequently
attained.

Finally, (11) transfers the same exact degree to residue \(k+1-j\). This
proves (19) throughout its stated range. \(\square\)

\subsection{8. Hypotheses, scope, and attribution}\label{verification-audit-and-provenance}

\textbf{Determinant convention.} The coefficient sequence is \(C^{2k}\),
the shift is zero, and the Hankel size is exactly \(kn+r\). The sign is
\((-1)^{n\binom{k}{2}}\). Theorem 1 treats both specified residues and
all \(n\ge0\), with no exceptional \(k\ge5\). The case \(k=5\) is
included once, not assumed to consist of two different residues.

\textbf{Theorem hypotheses.} The base polynomials are monic and
orthonormal for the explicitly constructed functional, whose norms are
all 1. The modifying polynomial is monic of even degree. Its simple
nonzero roots are distinct from the confluent point 0, so (7) is
nonzero. Series division uses \(M=k(n+1)\ne0\) in the target range. The
only analytic input is the classical cotangent expansion; it is used
inside its disk of convergence to identify coefficients, and then as a
formal coefficient identity. Every sum in the strict-positivity argument
converges.

\textbf{Exactness.} The positive coefficient in (17), rather than finite
interpolation or a proposed degree bound, establishes exact degree. The
algebraic negative-index extension is used solely for the proved
polynomial reflection identity. No negative-size Hankel determinant is
used.

\textbf{Scope.} Theorem 1 covers the stated residue-three case.
Theorem 7.1 further proves the full stated \(j\)-range in {[}C23,
equation (19){]}. The argument does not assert formulas for arbitrary
shifts, odd Catalan powers, or Narayana-parameter deformations.

\textbf{Method attribution.} The moment-functional construction,
determinant reduction, polynomiality, and reflection form the preliminary
framework and are included with proofs. The classical moment-modification
identity is attributed to {[}K21{]} in Section 2. The higher-rank quotient
minor and the strictly positive moment determinant supply the uniform
exact-degree argument. The supplementary finite certificate is not a
premise of this proof.

\textbf{Literature scope.} {[}C23{]} and {[}WX19{]} identify the target
and known finite range. {[}C18{]} supplies the neighboring Narayana
statements. {[}K21{]} gives the classical determinant transformation.
{[}H26{]} treats shifts \(-m,1-m,2-m,3-m\), which do not give the
present unshifted target for \(m=k\ge5\). The additional 2026 source
{[}CS26, Theorem 1.4{]}, specialized to its parameters
\(k_{\rm CS}=k-1,m=k-1,c=2\), gives the neighboring residue-0 and
residue-1 formulas; it does not state the residue-3 degree theorem. No
pre-existing uniform proof of the exact residue-degree statement was
identified in these checked sources. This is a limited literature check,
not a worldwide priority claim.

\paragraph{AI assistance.} This work was prepared with AI assistance.
No independent human peer-review or proof-assistant certification is
claimed.

\section{References}\label{citation}
\begingroup\small\raggedright
\setlength{\parskip}{3pt}

\textbf{{[}C23{]}} Johann Cigler, \emph{Some experimental observations
about Hankel determinants of convolution powers of Catalan numbers},
arXiv:2308.07642v2 (2023). Equations (1)--(2), page 1; equation (16) and
Conjecture 6, equation (19), page 4. This is the exact source of the
target and its range.

\textbf{{[}WX19{]}} Ying Wang and Guoce Xin, \emph{Hankel determinants
for convolution powers of Catalan numbers}, arXiv:1811.00248v2;
\emph{Discrete Mathematics} 342 (2019), 2694--2716. Conjecture 10, page
3, and Theorem 11, page 4. Only the specializations in the range
specified by {[}C23{]} are used here.

\textbf{{[}C18{]}} Johann Cigler, \emph{Catalan numbers, Hankel
determinants and Fibonacci polynomials}, arXiv:1801.05608v3 (2018).
Conjecture 7.7, equation (7.13), page 30. This concerns a neighboring problem, not
the residue-three statement.

\textbf{{[}K21{]}} Christian Krattenthaler, \emph{Hankel determinants of
linear combinations of moments of orthogonal polynomials, II},
arXiv:2101.04225v5 (2021). Theorem 1 and Proposition 5. DOI:
10.1007/s11139-021-00514-8.

\textbf{{[}H26{]}} Guo-Niu Han, \emph{Hankel Determinants for
Convolution Powers of Narayana Polynomials}, arXiv:2512.12603v1;
\emph{Electronic Journal of Combinatorics} 33(2) (2026), P2.42. Theorems
4--7. DOI: 10.37236/14995.

\textbf{{[}CS26{]}} Shane Chern and Wenle Shi, \emph{Hankel determinants
of Catalan-like sequences}, arXiv:2608.27208v1 (2026). Theorem 1.4, page
3.

\textbf{{[}DLMF{]}} NIST Digital Library of Mathematical Functions,
equations 4.22.3 (cotangent partial fractions) and 4.19.6 (cotangent
Laurent coefficients).

\endgroup
\end{document}
